\subsection{Lax squares and mates}
\label{subsec:mates-mates}

\begin{defi}
\label{def:lax-square}
\gutentag{000E}%
A \emph{lax square} in $\A$ is a diagram
\begin{equation*}
\begin{tikzcd}
A \ar[r, "f"] \ar[d, "g"'] & B \ar[d, "h"] \ar[dl, Rightarrow, shorten=3mm,
"\alpha"'] \\
C \ar[r, "k"'] & D
\end{tikzcd}
\end{equation*}
with a $2$-morphism
$\alpha\colon hf \Rightarrow kg$
from the top-right to the bottom-left composite.
It is a \emph{commutative square} if $\alpha$ is invertible.
\end{defi}

\gutentag{000F}%
The \emph{transpose} of a commutative square,
its reflection in the diagonal through $A$ and $D$,
is the commutative square with $g$, $k$ horizontal, $f$,
$h$ vertical and $2$-morphism
$\alpha^{-1}\colon kg \Rightarrow hf$.
So a commutative square has two underlying lax squares.

\gutentag{009C}%
\textbf{Important Convention:} we speak of the mates,
or of the Beck--Chevalley condition,
of a commutative square only after fixing one of its two underlying lax squares.
This way we can avoid saying ``horizontal/vertical left/right mate".
In particular, it could happen that we draw a lax square and take its mate,
while the lax square is in fact commutative.
Drawing a square as a lax square by no means implies that the square is not
commutative.

\begin{constru}[mates]
\label{constr:mates}
\gutentag{0010}%
Let $\alpha$ be the lax square of \cref{def:lax-square}.
\begin{enumerate}
\item If $f^{L} \dashv f$ and $k^{L} \dashv k$,
the \emph{left mate} of $\alpha$ is the lax square
\begin{equation*}
\begin{tikzcd}
B \ar[r, "h"] \ar[d, "f^{L}"'] & D \ar[d, "k^{L}"] \ar[dl, Rightarrow,
shorten=3mm, "\lmate{\alpha}"'] \\
A \ar[r, "g"'] & C
\end{tikzcd}
\qquad
\lmate{\alpha}\colon k^{L}h \xRightarrow{k^{L}h\eta_{f}} k^{L}hff^{L}
\xRightarrow{k^{L}\alpha f^{L}} k^{L}kgf^{L} \xRightarrow{\varepsilon_{k}gf^{L}}
gf^{L} \ .
\end{equation*}
\item If $g \dashv g^{R}$ and $h \dashv h^{R}$,
the \emph{right mate} of $\alpha$ is the lax square
\begin{equation*}
\begin{tikzcd}
C \ar[r, "g^{R}"] \ar[d, "k"'] & A \ar[d, "f"] \ar[dl, Rightarrow, shorten=3mm,
"\rmate{\alpha}"'] \\
D \ar[r, "h^{R}"'] & B
\end{tikzcd}
\qquad
\rmate{\alpha}\colon fg^{R} \xRightarrow{\eta_{h}fg^{R}} h^{R}hfg^{R}
\xRightarrow{h^{R}\alpha g^{R}} h^{R}kgg^{R} \xRightarrow{h^{R}k\varepsilon_{g}}
h^{R}k \ .
\end{equation*}
\end{enumerate}
Thus the left mate replaces the horizontal $1$-morphisms by their left adjoints
(rotating the square by $90^{\circ}$ to the left),
and the right mate replaces the vertical $1$-morphisms by their right adjoints
(rotating it to the right). Here $\eta$ and $\varepsilon$ denote units and
counits.
\end{constru}

\begin{defi}
\label{def:BC}
\gutentag{0011}%
A lax square is \emph{left Beck--Chevalley} if $f$,
$k$ admit left adjoints and $\lmate{\alpha}$ is invertible,
and \emph{right Beck--Chevalley} if $g$,
$h$ admit right adjoints and $\rmate{\alpha}$ is invertible.
\end{defi}

\begin{lem}[the two mates are conjugate]
\label{lem:mates-conjugate}
\gutentag{009D}%
Let $\alpha$ be the lax square of \cref{def:lax-square}, with $f^{L} \dashv f$,
$k^{L} \dashv k$, $g \dashv g^{R}$ and $h \dashv h^{R}$. Then
$\rmate{\alpha}\colon fg^{R} \Rightarrow h^{R}k$
is the right mate of
$\lmate{\alpha}\colon k^{L}h \Rightarrow gf^{L}$,
viewed as a lax square with identity horizontal edges and vertical edges $gf^{L}
\dashv fg^{R}$ and $k^{L}h \dashv h^{R}k$. In particular,
$\alpha$ is left Beck--Chevalley if and only if it is right Beck--Chevalley.
\end{lem}

\begin{proof}
\gutentag{009E}%
The right mate of $\lmate{\alpha}$ is the composite
\begin{equation*}
fg^{R} \xRightarrow{\ \eta\ } h^{R}k\,k^{L}h\,fg^{R}
\xRightarrow{\ h^{R}k\,\lmate{\alpha}\,fg^{R}\ } h^{R}k\,gf^{L}\,fg^{R}
\xRightarrow{\ \varepsilon\ } h^{R}k \ ,
\end{equation*}
where
$\eta = h^{R}\eta_{k}h \circ \eta_{h}$
is the unit of $k^{L}h \dashv h^{R}k$ and
$\varepsilon = \varepsilon_{g} \circ g\varepsilon_{f}g^{R}$
the counit of $gf^{L} \dashv fg^{R}$. Insert the formula of \cref{constr:mates}
for $\lmate{\alpha}$. By interchange in $\htwo\A$,
$\eta_{k}$ and $\varepsilon_{k}$ cancel by the triangle identity
$k\varepsilon_{k} \circ \eta_{k}k = \mathrm{id}_{k}$,
and $\eta_{f}$ and $\varepsilon_{f}$ by
$f\varepsilon_{f} \circ \eta_{f}f = \mathrm{id}_{f}$.
What remains is
$h^{R}k\varepsilon_{g} \circ h^{R}\alpha g^{R} \circ \eta_{h}fg^{R} =
\rmate{\alpha}$.
For the last claim: by \cref{lem:mate-correspondence}(1) and (3),
taking mates of squares with identity horizontal edges is a bijection compatible
with composition and identities, so it preserves and reflects invertibility.
\end{proof}

\begin{rmk}
\label{rmk:double-BC}
\gutentag{009F}%
The lemma does not extend to the \emph{double Beck--Chevalley map} of
Cnossen--Lenz--Linskens \cite[Def.~2.1(4)]{cnossenlenzlinskens2025universal}.
If $\alpha$ is left Beck--Chevalley, replacing $f$ and $k$ by their left
adjoints gives the \emph{mirror} of $\alpha$, the commutative square
\begin{equation*}
\begin{tikzcd}
B \ar[r, "f^{L}"] \ar[d, "h"'] & A \ar[d, "g"] \ar[dl, Rightarrow, shorten=3mm,
"{\lmate{\alpha}}^{-1}"'] \\
D \ar[r, "k^{L}"'] & C
\end{tikzcd}
\end{equation*}
whose right mate
$\delta\colon f^{L}h^{R} \Rightarrow g^{R}k^{L}$
is the double Beck--Chevalley map. Following
\cite{cnossenlenzlinskens2025universal},
$\alpha$ is \emph{biadjointable} if it is left Beck--Chevalley and $\delta$ is
invertible. This is a further condition: in $\Cat$, let
$u\colon \Delta^{1} \to \Delta^{2}$
be the inclusion of $\{0,2\}$, which is fully faithful with both adjoints,
and let $\alpha$ be the identity square with
$f = g = \mathrm{id}_{\Delta^{1}}$
and $h = k = u$. Its mates are the counit
$u^{L}u \simeq \mathrm{id}$
and the unit
$\mathrm{id} \simeq u^{R}u$,
both invertible, while
$\delta\colon u^{R} \Rightarrow u^{L}$
is at $1 \in \Delta^{2}$ the morphism $0 \to 1$ of $\Delta^{1}$.
\end{rmk}

\begin{lem}[biadjointable squares]
\label{lem:biadjointable}
Let $\E$ be an $(\infty,2)$-category, and draw
$\Delta^{1} \times \Delta^{1}$
with the first factor horizontal.
\begin{enumerate}
\item
$(\rho \times \ell)^{*}\colon \func(\Adj \times \Adj,\E) \to
\func(\Delta^{1} \times \Delta^{1},\E)$
is a locally full inclusion. Its objects are the commutative squares $\alpha$ as
in \cref{def:lax-square} with $f$, $k$ right adjoints and $g$, $h$ left adjoints
that are biadjointable; equivalently, that are right Beck--Chevalley and such
that replacing $g$, $h$ by $g^{R}$, $h^{R}$ and filling with
${\rmate{\alpha}}^{-1}$
gives a left Beck--Chevalley square. Its $1$-morphisms are the transformations
whose naturality squares at $f$ and $k$ are left, and at $g$ and $h$ right,
Beck--Chevalley.
\item
$(\rho \times \ell)^{*}\colon \func(\Refl \times \Corefl,\E) \to
\func(\Delta^{1} \times \Delta^{1},\E)$
is a locally full inclusion, with objects and $1$-morphisms as in (1) except
that $f$, $k$ are reflective and $g$, $h$ coreflective. Here $\Refl$ and
$\Corefl$ are the walking reflection and coreflection, which we will define in
\cref{subsec:mates-morphisms}.
\end{enumerate}
\end{lem}

\begin{proof}
\gutentag{00A1}%
The proof uses \cref{lem:mates} and \cref{lem:mates-refl},
which we will prove in \cref{subsec:mates-morphisms} without using this lemma.
(1) is \cite[Prop.~2.12(3)]{cnossenlenzlinskens2025universal};
we recall the argument.
Factor
$(\rho \times \ell)^{*}$
as restriction along
$\mathrm{id} \times \ell$
followed by restriction along
$\rho \times \mathrm{id}$,
i.e.\ as
\begin{equation*}
\func\bigl(\Adj,\func(\Adj,\E)\bigr) \xrightarrow{\ (\ell^{*})_{*}\ }
\func\bigl(\Adj,\func(\Delta^{1},\E)\bigr) \simeq
\func\bigl(\Delta^{1},\func(\Adj,\E)\bigr) \xrightarrow{\ (\rho^{*})_{*}\ }
\func\bigl(\Delta^{1},\func(\Delta^{1},\E)\bigr) \ ,
\end{equation*}
with $\Delta^{1}$ the vertical factor. Both maps are locally full inclusions by
\cref{lem:mates} and \cref{subsec:prelim-conventions},
so we compute images with
\cite[Cor.~2.6.6]{abellangagnahaugseng2024straightening}.
By \cref{lem:mates}(2),
the objects in the image of $(\rho^{*})_{*}$ are the squares with $f$,
$k$ right adjoints that are left Beck--Chevalley.
Such an $\alpha$ lifts to a transformation between the adjunctions $f^{L} \dashv
f$ and $k^{L} \dashv k$, whose naturality square at $r$ is $\alpha$ and at $l$
is the mirror of $\alpha$, as in the proof of \cref{lem:mates}
\cite[Cor.~4.9]{haugseng2021laxmonads}. Viewed as
$G\colon \Adj \to \func(\Delta^{1},\E)$,
it lies in the image of $(\ell^{*})_{*}$ if and only if $G(\pm)$,
i.e.\ $g$ and $h$, are left adjoints and $G(r)$, $G(l)$,
i.e.\ $\alpha$ and its mirror, are right Beck--Chevalley
(\cref{lem:mates}(1));
here we use that the $1$-morphisms of $\Adj$ are composites of $l$ and $r$ and
that Beck--Chevalley squares are closed under pasting.
As $\alpha$ is right Beck--Chevalley if and only if it is left Beck--Chevalley
(\cref{lem:mates-conjugate}), this is the first description of the objects.
Factoring through
$\func(\Delta^{1} \times \Adj,\E)$
instead gives the second in the same way.
A $1$-morphism between squares in the image is a functor from
$\Delta^{1} \times \Delta^{1}$,
the first factor being the direction of the transformation and the second the
vertical one, to $\func(\Delta^{1},\E)$.
It lies in the image of $(\rho^{*})_{*}$ if and only if its naturality squares
at $f$ and $k$ are left Beck--Chevalley, as the other $1$-morphisms of
$\Delta^{1} \times \Delta^{1}$
go to the given squares or to composites;
likewise its lift lies in the image of $(\ell^{*})_{*}$ if and only if its
naturality squares at $g$ and $h$ are right Beck--Chevalley.
(2) The argument for $\Adj$ goes through verbatim,
with \cref{lem:mates-refl}(2) in place of \cref{lem:mates}.
\end{proof}

\begin{lem}[mate correspondence]
\label{lem:mate-correspondence}
\gutentag{0012}%
In the situation of \cref{constr:mates}:
\begin{enumerate}
\item
$\alpha \mapsto \lmate{\alpha}$
and
$\alpha \mapsto \rmate{\alpha}$
are equivalences of spaces
\begin{equation*}
\mapp_{\Hom(B,C)}(k^{L}h,gf^{L}) \simeq \mapp_{\Hom(A,D)}(hf,kg) \simeq
\mapp_{\Hom(C,B)}(fg^{R},h^{R}k) \ .
\end{equation*}
\item They are mutually inverse: the right mate of $\lmate{\alpha}$ (with
respect to $f^{L} \dashv f$, $k^{L} \dashv k$) and the left mate of
$\rmate{\alpha}$ (with respect to $g \dashv g^{R}$,
$h \dashv h^{R}$) are $\alpha$.
\item Mates are compatible with horizontal and vertical pasting and with
identity squares.
\end{enumerate}
\end{lem}

\begin{proof}
\gutentag{0013}%
(1) The adjunctions
$(k^{L})_{\natural} \dashv k_{\natural}$
and
$f^{\natural} \dashv (f^{L})^{\natural}$
of \eqref{eq:hom-adjunctions} give
\begin{equation*}
\mapp(hf,kg) \simeq \mapp(k^{L}hf,g) \simeq \mapp(k^{L}h,gf^{L}) \ ,
\end{equation*}
and unwinding units and counits shows that the composite is
$\alpha \mapsto \lmate{\alpha}$.
Similarly
$h_{\natural} \dashv (h^{R})_{\natural}$
and
$(g^{R})^{\natural} \dashv g^{\natural}$
give
\begin{equation*}
\mapp(hf,kg) \simeq \mapp(f,h^{R}kg) \simeq \mapp(fg^{R},h^{R}k) \ .
\end{equation*}
(2) and (3) follow from the triangle identities and interchange in $\htwo\A$ as
in \cite{kellystreet1974review}.
\end{proof}

\begin{rmk}
\label{rmk:mates-memo}
\gutentag{0014}%
If $\alpha$ is commutative with transpose $\alpha^{\mathrm{T}}$,
the left mate of $\alpha^{\mathrm{T}}$ replaces $g$,
$h$ by their left adjoints and its right mate replaces $f$,
$k$ by their right adjoints; by our convention these are mates of
$\alpha^{\mathrm{T}}$, not of $\alpha$.
\end{rmk}
